\documentclass{article}

\begin{document}
    \begin{center}
        \Large \textbf{Challenge 2: Terraform und NixOS} \\ [6pt]
        \normalsize
        Autor: Maximilian Jakob Maag B.Sc. \\
        Version: 1.0
    \end{center}
    \section{Intro}
    Terraform ist hervorragend dafür geeignet, um Infrastruktur zu provisionieren. Allerdings ist es meist unproduktiv, ein Element einer Infrastruktur zu zerstören und neu anzulegen,
    um eine kleine Änderung vorzunehmen. Terraform lässt sich auf verschiedene Wege hervorragend mit anderen Technologien kombinieren.
    NixOS ist eine Linux-Distribution, die durch eine Konfigurationsdatei beschrieben wird.
    Der Sollzustand, z.B. immer die neueste Version eines Pakets zu installieren oder Updates, lassen sich super automatisieren.

    \section{Aufgaben}
    Diese Challenge soll mehrere Welten miteinander verbinden. Terraform als Werkzeug für Provisionierung und NixOS für Wartungsaufgaben auf einem VPS.
    Terraform soll Aufgaben um das VPS Übernehmen, während das VPS selbst nicht immer zerstört und neu angelegt wird. Beispielsweise Firewalls oder Portfreigaben lassen sich super mit Terraform ändern.
    NixOS soll für die Wartung des VPS zuständig sein. Beispielsweise das Installieren von Paketen, das Einrichten von Diensten oder das Anlegen von Benutzern. \\ \\
    Diese Übung konzentriert sich auf Wartungsaufgaben nach dem Deployment mittels Terraform. NixOS ist eines von vielen Beispielen; Ansible ist ebenfalls eine passende Lösung.

    \subsection{Aufgabe 1}
    NixOS ist nicht bei allen Providern standardmäßig vorhanden.
    Deployen Sie ein VPS mit NixOS in einer Region Ihrer Wahl. Alle weiteren Aufgaben müssen ohne Terraform erfüllt werden.

    \subsection{Aufgabe 2}
    Erstellen Sie eine NixOS-Konfiguration, welche einen Benutzer hinzufügt. Außerdem soll ein Webserver installiert werden.
    Sie können die NixOS-Konfiguration auch von einer KI erzeugen lassen. \\ \\
    Verbinden Sie sich mittels SSH mit Ihrem deployten VPS und wenden Sie die NixOS-Konfiguration an. Überprüfen Sie die Anwendung, indem Sie den Webserver aufrufen.

    \subsection{Aufgabe 3}
    Verbinden Sie sich erneut mit Ihrem VPS. Führen Sie ein Rollback auf den Ausgangszustand von NixOS durch, ohne Terraform zu verwenden.

    \subsection{Aufgabe 4}
    Ändern Sie die Firewallregeln Ihres VPS mittels NixOS. Lassen Sie nur den Port 22 zu.
    
    \subsection{Aufgabe 5}
    Deployen Sie nun mittels Terraform eine Firewall, die ebenfalls nur port 22 offen lässt, ohne das VPS zu zerstören.
\end{document}